% GENERATED FILE. Edit canonical lessons, then run npm run generate:books.
% canonical-source-hash: fnv1a64:9ed1244e53a2a597
% canonical-lessons: ES-C330-camino, ES-C330-parada, ES-C330-mundo, ES-C330-vida, ES-C330-vez

\chapter{The Road and the World}
\label{ch:camino-y-vez}
\hlchaptermodality{330}

\begin{chapteropening}
\textbf{By the end of this chapter:} \emph{I can name el camino, la parada, el mundo, la vida and la vez, and say which of them is Celtic rather than Latin.}

Complete ES-C330-vez: follow a camino to a parada, widen out to el mundo and la vida, and open the noun inside otra vez.
\end{chapteropening}

\section[el camino]{el camino — a Celtic word that outlived the Romans}

\label{lesson:ES-C330-camino}



\begin{quote}
Say \emph{but rather}, then \emph{except}. (\emph{Sino}; \emph{salvo}.)
\end{quote}

\subsection*{What to know first}
\begin{itemize}
\raggedright
  \item el lugar — the spot a road runs between.
\end{itemize}

\begin{sounds}
\begin{itemize}
\raggedright
  \item \emph{ca-MI-no}, three beats with the weight in the middle. The single \emph{m} and single \emph{n} are both light taps; neither doubles. $\to$ pronunciation reference
\end{itemize}
\end{sounds}

\begin{cousinweb}
\textbf{el camino} — \textquotedblleft{}the road,\textquotedblright{} \textquotedblleft{}the way.\textquotedblright{}

Almost every Spanish noun you have met so far has come down from Latin. This one did not start there.

Before Rome reached the Iberian peninsula and Gaul, Celtic-speaking peoples lived across both, and they had a word something like \textbf{cammanos} for a track or a way. The Romans, who were famous road-builders and had perfectly good words of their own, borrowed the local one anyway. It shows up in Late Latin as \textbf{camminus}, and from there it became Spanish \textbf{camino}, French \textbf{chemin}, and Italian \textbf{cammino}.

Very little Celtic survives in Spanish — a few dozen words for landscape, carts and beer — so this is a genuine rarity. When you say \emph{camino} you are using a word that was in the ground before Latin arrived and refused to be replaced.

One collision to keep straight, because it looks like a match and is not. Latin also had \textbf{camīnus}, meaning a \emph{furnace} or \emph{hearth}, borrowed from Greek \textbf{kaminos}. That word gave English \textbf{chimney}. It has nothing to do with roads and nothing to do with this word. Two different sources landed on nearly the same spelling, and the resemblance is an accident.
\end{cousinweb}

\begin{grammarlens}[title={What you've built}]
The first of five nouns for the shape of a day: where you go, and how long things take.
\end{grammarlens}

\subsection*{Your turn}
Say these aloud:

\begin{itemize}
\raggedright
  \item \emph{el camino}
  \item \emph{un camino}
  \item \emph{el camino a casa}
\end{itemize}

\begin{tcolorbox}[breakable,colback=teal!4,colframe=teal!35!black,title={Before you move on}]
Which language did \emph{camino} come from before Latin? (Celtic — Gaulish.) And which lookalike Latin word gave English \emph{chimney}? (\emph{Camīnus}, a furnace.)
\end{tcolorbox}
\section[la parada]{la parada — the stop that used to mean \emph{get ready}}

\label{lesson:ES-C330-parada}



\begin{quote}
Say \emph{except}, then \emph{the road}. (\emph{Salvo}; \emph{el camino}.)
\end{quote}

\subsection*{What to know first}
\begin{itemize}
\raggedright
  \item ¿dónde está? — the question you ask about one of these.
\end{itemize}

\begin{sounds}
\begin{itemize}
\raggedright
  \item \emph{pa-RA-da}, three beats, the weight in the middle. The \emph{r} is a single quick tap of the tongue, and the \emph{d} between vowels is soft, close to the \emph{th} in \emph{this}. $\to$ pronunciation reference
\end{itemize}
\end{sounds}

\begin{cousinweb}
\textbf{la parada} — \textquotedblleft{}the stop,\textquotedblright{} the place a bus or a train halts.

It is built from the verb \textbf{parar}, \emph{to stop}, and \emph{parar} comes from Latin \textbf{parāre} — which did not mean \emph{to stop} at all. It meant \textbf{to make ready, to prepare}.

Follow the drift. To \emph{make ready} is to set a thing in position. To set it in position is to place it. To place it is to make it stand still. Spanish walked that path over centuries and arrived at \emph{stopping}; the other Romance languages mostly did not, which is why this meaning surprises people who know French or Italian.

English took \emph{parāre} over and over and kept the original sense every time. To \textbf{prepare} is to make ready beforehand. To \textbf{repair} is to make ready again. An \textbf{apparatus} is a set of things made ready. To \textbf{parade} is to make a show ready. A \textbf{parasol} is made ready \emph{against} the sun, and a \textbf{parachute} ready against a fall — that \emph{para-} is the Romance \emph{parare}, and not the Greek \emph{para-} meaning \emph{beside} that sits in \emph{parallel} and \emph{parable}.

So English has a dozen words insisting the root means \emph{prepare}, and Spanish has the bus stop insisting it means \emph{halt}. Both are honest descendants.
\end{cousinweb}

\begin{grammarlens}[title={What you've built}]
Two of five. A road, and the places on it where you get off.
\end{grammarlens}

\subsection*{Your turn}
Say these aloud:

\begin{itemize}
\raggedright
  \item \emph{la parada}
  \item \emph{una parada}
  \item \emph{¿dónde está la parada?}
\end{itemize}

\begin{tcolorbox}[breakable,colback=teal!4,colframe=teal!35!black,title={Before you move on}]
What did Latin \emph{parāre} mean? (To make ready, to prepare.) And which sense did Spanish end up with? (Stopping.)
\end{tcolorbox}
\section[el mundo]{el mundo — the world, which used to mean \emph{tidy}}

\label{lesson:ES-C330-mundo}



\begin{quote}
Say \emph{the road}, then \emph{the stop}. (\emph{El camino}; \emph{la parada}.)
\end{quote}

\subsection*{What to know first}
\begin{itemize}
\raggedright
  \item el país — the widest circle you had until now.
\end{itemize}

\begin{sounds}
\begin{itemize}
\raggedright
  \item \emph{MUN-do}, two beats with the weight on the first. The \emph{n} before \emph{d} is made with the tongue already at the teeth, so the two run together. $\to$ pronunciation reference
\end{itemize}
\end{sounds}

\begin{cousinweb}
\textbf{el mundo} — \textquotedblleft{}the world.\textquotedblright{}

Latin \textbf{mundus} started life as an adjective, and it meant \textbf{clean, neat, orderly, elegant}. A \emph{mundus} room was a tidy one. The opposite, \textbf{immundus}, meant filthy — and Spanish still has \textbf{inmundo} for something foul.

So how did \emph{tidy} become \emph{the world}?

Roman writers needed a Latin word for the Greek \textbf{kósmos}. And \emph{kósmos} had made this exact journey already: it first meant \emph{order, arrangement, adornment}, and Greek philosophers stretched it to name the universe — because what impressed them about the universe was that it was \textbf{arranged}, not chaotic. To call the world a \emph{kosmos} was to make a claim about it.

Latin needed a word that already meant \emph{orderly} to translate a word that meant \emph{orderly}. \textbf{Mundus} was the natural pick, and it inherited the second meaning along with the first.

Both halves of the Greek word survive in English, and they look unrelated until you know this. The \textbf{cosmos} is the ordered universe. \textbf{Cosmetics} are ordering and adornment applied to a face. Same word, two sizes.

And \emph{mundus} survives in English too, in \textbf{mundane} — which now means \emph{dull and everyday}, having slid from \emph{of the world} to \emph{worldly} to \emph{ordinary}. A word that once meant \emph{elegant} now means \emph{boring}.
\end{cousinweb}

\begin{grammarlens}[title={What you've built}]
Three of five. You can now speak about somewhere as small as a bus stop and as large as everything there is.
\end{grammarlens}

\subsection*{Your turn}
Say these aloud:

\begin{itemize}
\raggedright
  \item \emph{el mundo}
  \item \emph{todo el mundo}
  \item \emph{el mundo es grande}
\end{itemize}

\begin{tcolorbox}[breakable,colback=teal!4,colframe=teal!35!black,title={Before you move on}]
What did \emph{mundus} mean as an adjective? (Clean, orderly.) And which Greek word made the same jump from order to world? (\emph{Kósmos}.)
\end{tcolorbox}
\section[la vida]{la vida — the root that vitamins were named for}

\label{lesson:ES-C330-vida}



\begin{quote}
Say \emph{the stop}, then \emph{the world}. (\emph{La parada}; \emph{el mundo}.)
\end{quote}

\subsection*{What to know first}
\begin{itemize}
\raggedright
  \item pasar a mejor vida — an idiom this word already carries.
\end{itemize}

\begin{sounds}
\begin{itemize}
\raggedright
  \item \emph{VI-da}, two beats, the weight on the first. The \emph{v} is soft, made without closing the lips, and the \emph{d} between vowels is soft too. $\to$ pronunciation reference
\end{itemize}
\end{sounds}

\begin{cousinweb}
\textbf{la vida} — \textquotedblleft{}life.\textquotedblright{}

Latin \textbf{vīta}. Spanish softened the \emph{t} between vowels into a \emph{d}, as it did in hundreds of words, and \emph{vīta} became \textbf{vida}.

English holds the root everywhere. Something \textbf{vital} is necessary for life. To \textbf{revive} is to bring back to life, to \textbf{survive} is to live past something, and a plan that is \textbf{viable} is one that can live.

The best story belongs to \textbf{vitamin}. In 1912 the Polish biochemist Casimir Funk found a substance in rice husks that cured beriberi, and concluded that such life-critical compounds were \textbf{amines}, a family of nitrogen compounds. He named them \textbf{vitamines} — \emph{vita} plus \emph{amine}, life-amines.

He was half right. They were essential to life, but most of them were not amines at all. Rather than rename everything, chemists quietly dropped the final \emph{e}, and \textbf{vitamine} became \textbf{vitamin} — a word whose spelling still records a mistake that was corrected by deleting one letter.

Spanish uses \emph{vida} in more places than English uses \emph{life}. \textbf{Así es la vida} is \emph{that is how life is}, and you now have every word in it.
\end{cousinweb}

\begin{grammarlens}[title={What you've built}]
Four of five. Road, stop, world, life — the frame around everything else you have learned to say.
\end{grammarlens}

\subsection*{Your turn}
Say these aloud:

\begin{itemize}
\raggedright
  \item \emph{la vida}
  \item \emph{así es la vida}
  \item \emph{mi vida}
\end{itemize}

\begin{tcolorbox}[breakable,colback=teal!4,colframe=teal!35!black,title={Before you move on}]
What did Latin \emph{vīta} mean? (Life.) And why did \emph{vitamine} lose its final \emph{e}? (Most of them turned out not to be amines.)
\end{tcolorbox}
\section[la vez]{la vez — the noun inside \emph{otra vez} and \emph{tal vez}}

\label{lesson:ES-C330-vez}



\begin{quote}
Say \emph{the world}, then \emph{life}. (\emph{El mundo}; \emph{la vida}.)
\end{quote}

\subsection*{What to know first}
\begin{itemize}
\raggedright
  \item otra vez — the repair phrase you already use.
  \item tal vez — the hedge, whose \emph{tal} you met two chapters ago.
\end{itemize}

\begin{sounds}
\begin{itemize}
\raggedright
  \item \emph{VES}, one beat. The \emph{v} is soft. The final \emph{z} is a plain hissed \emph{s} across the Americas, and a \emph{th} sound in much of Spain. $\to$ pronunciation reference
\end{itemize}
\end{sounds}

\begin{cousinweb}
\textbf{la vez} — \textquotedblleft{}a time,\textquotedblright{} in the sense of \emph{an occasion} or \emph{a turn}.

Spanish, like English, keeps two different words apart here. \textbf{El tiempo} is time as a substance, the stuff that passes. \textbf{La vez} is a countable occurrence: \emph{one time, three times, the last time}.

Latin \textbf{vicem} was the source, and it meant \emph{a turn, a change, a place taken in alternation}. The core idea is \textbf{substitution}: my turn is the slot I occupy when it is not yours.

You have been saying it for a long time already. \textbf{Otra vez} is \emph{another turn} — which is why it means \emph{again}. And \textbf{tal vez} is \emph{such a turn}, \emph{a turn of that kind}, which is how Spanish arrives at \emph{maybe}. Now that you have \emph{tal} from the chapter before this one, both words in that phrase are ones you can take apart.

English took the substitution sense and made a career of it. A \textbf{vice-president} holds office \emph{in the turn of} the president. A \textbf{viceroy} ruled in the king's place. A \textbf{vicar} stands in for a higher authority. And something learned \textbf{vicariously} is learned through somebody else's experience standing in for your own.

There is one more, and it is everywhere: \textbf{vice versa}, \emph{the turn having been turned} — the positions swapped.
\end{cousinweb}

\begin{grammarlens}[title={What you've built}]
Five nouns: \emph{el camino}, \emph{la parada}, \emph{el mundo}, \emph{la vida}, \emph{la vez}. Two of them were hiding inside phrases you have used for hundreds of lessons.
\end{grammarlens}

\subsection*{Your turn}
Say these aloud:

\begin{itemize}
\raggedright
  \item \emph{la vez}
  \item \emph{otra vez}
  \item \emph{tal vez}
\end{itemize}

\begin{tcolorbox}[breakable,colback=teal!4,colframe=teal!35!black,title={Before you move on}]
What did Latin \emph{vicem} mean? (A turn, a change of place.) And what does \emph{otra vez} say word for word? (Another turn.)
\end{tcolorbox}
